\centerline{\bf \Huge Вступительный экзамен. Практика}

\problem{\textit{\textbf{1}}} а.\textbf{(3)} Переведите в десятичную дробь число $\dfrac{2026}{23}$.

б.\textbf{(3)} Переведите в обыкновенную дробь число $2504,(2026)$.


\problem{\textit{\textbf{2}}}\textbf{(9)} Найдите НОД и НОК тройки следующих чисел: 2362500,3207600,3744000.

\problem{\textit{\textbf{3}}}\textbf{(12)} Лайт написал на доске число, которое делится на 24. Затем стёр в нём цифры, и на доске осталась следующая запись 2504?2026?? (? - стёртые цифры). Какое число изначально мог написать на доске Лайт? Перечислите ВСЕ варианты и докажите, что других быть не может.

\problem{\textit{\textbf{4}}}\textbf{(15)} Чему равно следующее выражение: 

$\dfrac{1572864}{1048576}\cdot\Bigr(\dfrac{(1,238+2,762) \cdot \dfrac{11}{110}}{(36,487-34,237): 2,8125}+\dfrac{(4,36-1,16) \cdot 0,3125}{\dfrac{9}{45} \cdot(47,8-45,55): 0,225}\Bigl)\cdot$ 

$\cdot \Bigr(39,86-\dfrac{8 \dfrac{4}{5}}{6,934+38 \dfrac{13}{40}: 2 \dfrac{1}{2} \cdot 0,2}-17\Bigl)\cdot\Bigr(7,5^2-2,5^2\Bigl)\cdot\Bigr(\dfrac{20\dfrac{8}{15}\cdot7,5-54,6:\dfrac{2}{5}}{3\dfrac{13}{21}\cdot8,4-34,4:14\dfrac{1}{3}}\Bigl)\cdot12,8--3188=?$

\problem{\textit{\textbf{5}}} Решите системы линейных алгебраических уравнений (СЛАУ):

а.\textbf{(3)} $ \begin{cases}
   x + 2y = 5
   \\
   10x + 22y = 52
 \end{cases}$


б.\textbf{(9)}$\begin{cases}
   x + 2y + 3z = 12
   \\
   2x - y + z = -1
   \\
   3x + 7y - 2z = 5
   
 \end{cases}$

\problem{\textit{\textbf{6}}}а.\textbf{(6)} Прямая проходит через точку I(6;7) и точку T(-3;8). Найдите уравнение прямой и постройте её.

б.\textbf{(6)} Прямая $y=2x-11$ пересекается с прямой $y=-3x-2$ в точке $R$. Найдите ТОЧНЫЕ координаты точки $R$ и постройте графики двух прямых в одной декартовой системе координат.

\newpage

\problem{\textit{\textbf{7}}}\textbf{(12)} АА купил своим ученикам 20 кг сникерсов и 30 кг киндеров, заплатив за всё 210 тысяч рублей. Если сникерсы подорожают на 30\%, а киндеры подешевеют на 10\%, то за такую же массу соответствующих сладостей АА заплатит 213 тысяч рублей. Сколько рублей стоит 100г каждой сладости?


\problem{\textit{\textbf{8}}}\textbf{(6)} Найдите пересечение и объединение множеств $I$ и $T$: (\{-1917\} - одно число -1917)

$I \in (-2026;-2002] \cup [-1971;-1953) \cup [-1945,0905;-1941] \cup \{-1917\} \cup (-1162;988]$ 

$T \in (-2012;-2002) \cup (-2000;-1941) \cup (-1924;-1914] \cup [-1161;2022)$ 



\problem{\textit{\textbf{9}}}\textbf{(15)} Вычислите, раскрыв скобки: $(a-b)^{15} - (a+b)^{15}=?$

\problem{\textit{\textbf{10}}}\textbf{(9)} Известно, что площадь прямоугольника с целочисленными длинами сторон равна 432, а периметр больше 86 и меньше 437. Найдите всевозможные варианты длин сторон прямоугольника и докажите, что других быть не может.

\problem{\textit{\textbf{11}}} Решите уравнения:

а,б,в.\textbf{(по 3)} $x^2-7x+13=0$ , $3x^{2}-30x+75=0$, $-6x^{2}+x+7=0$

г.\textbf{(9)} $(2x)^{4}+2(2x)^{2}-63=0$

д.\textbf{(18)} $(x^2+2x+1-1)^2-(x+1)^2-55=0$

\problem{\textit{\textbf{12}}} Решите линейные диофантовы уравнения с двумя целыми неизвестными:

а.\textbf{(12)} $25x+4y=-10$

б.\textbf{(18)} $222156x-248292y=209088$


\problem{\textit{\textbf{13}}}\textbf{(12)} Чему равны целые числа - $L$ и $N$, если верно следующее равенство: $\dfrac{L}{7}-\dfrac{N}{11}=\dfrac{3}{77}$



\problem{\textit{\textbf{14}}} Сравните числа (выбор без объяснений 0 баллов): 

а\textbf{(6)} $\sqrt{37}$ vs $\sqrt{77}-\sqrt{7}$

б.\textbf{(15)} $\sqrt{\sqrt{\sqrt{3+67\sqrt{2}}}}$ vs $\sqrt{\sqrt{2}+\sqrt{3}}$


\problem{\textit{\textbf{15}}} Чему равны пределы следующих числовых последовательностей:

а.\textbf{(9)}$\lim_{n\to \infty} \dfrac{25n^2-n+1000-22n^3+2025n^6}{7-5n^5+19n^4+45n^6-9n}=?$

б.\textbf{(21)}$\lim_{n\to \infty} \dfrac{(1+2+3+...+(n-1)+n)\cdot(4+8+12+...+4(n-1)+4n)}{(2+4+6+...+2(n-1)+2n)\cdot (n+1)^2}=?$
